\section{Basics}
\label{sec:basics}

In this section, we introduce the denotational semantics of simple inductive types through a worked example, and then show how the construction generalises to containers and W-types. We then introduce
the category of containers.

Informally, the denotational semantics of inductive types is a correspondence between strictly-positive functors and their least fixed points, the latter being representable as strictly-positive inductive types. For example, the $\mathsf{Maybe}$ functor, defined as the operation $X \mapsto \mathbf{1} + X$ on \AgdaPrimitive{Type}, has the natural numbers \AgdaDatatype{ℕ} as its least fixed point.

Formally, given any \AgdaBound{A : }\AgdaPrimitive{Type}, and any F-algebra \AgdaBound{α :} \AgdaPrimitive{Maybe }\AgdaBound{A → A}, there is an unique folding function \AgdaFunction{foldℕ}\AgdaSpace{}%
\AgdaSymbol{:}\AgdaSpace{}%
\AgdaSymbol{(}\AgdaFunction{Maybe}\AgdaSpace{}%
\AgdaGeneralizable{A}\AgdaSpace{}%
\AgdaSymbol{→}\AgdaSpace{}%
\AgdaGeneralizable{A}\AgdaSymbol{)}\AgdaSpace{}%
\AgdaSymbol{→}\AgdaSpace{}%
\AgdaDatatype{ℕ}\AgdaSpace{}%
\AgdaSymbol{→}\AgdaSpace{}%
\AgdaGeneralizable{A}, such makes the following diagram commutes:

\[
\begin{tikzcd}[column sep=large, row sep=large]
  \mathsf{Maybe}\,\mathbb{N}
    \arrow[r, "\mathsf{Maybe}\,(\mathsf{fold}\,\alpha)"]
    \arrow[d, "{[\mathsf{zero},\,\mathsf{suc}]}"'] &
  \mathsf{Maybe}\,A
    \arrow[d, "\alpha"] \\
  \mathbb{N}
    \arrow[r, dashed, "\mathsf{fold}\,\alpha"'] &
  A
\end{tikzcd}
\]

To formalise this explicitly, it suffices to construct two functions: \AgdaFunction{commuteℕ}, which witnesses that \AgdaFunction{foldℕ} is a homomorphism of $\mathsf{Maybe}$-algebras, and \AgdaFunction{uniquefoldℕ}, which proves the uniqueness. Together, these establish that ℕ is the initial $\mathsf{Maybe}$-algebra.

\paragraph*{Containers}
\begin{itemize}
  \item intuition
  \item defition
  \item extension functor
\end{itemize}
